% Generated draft: requires mathematical and source review.
\documentclass[11pt]{article}
\usepackage[margin=25mm]{geometry}
\usepackage{fontspec}
\setmainfont{Latin Modern Roman}
\usepackage{amsmath,amssymb,mathtools,unicode-math}
\setmathfont{Latin Modern Math}
\usepackage{enumitem,needspace,xurl,hyperref}
\setlength{\parindent}{0pt}
\setlength{\parskip}{5pt}
\setlength{\emergencystretch}{3em}
\newcommand{\N}{\mathbb{N}}
\newcommand{\Z}{\mathbb{Z}}
\newcommand{\Q}{\mathbb{Q}}
\newcommand{\R}{\mathbb{R}}
\newcommand{\C}{\mathbb{C}}
\newcommand{\F}{\mathbb{F}}
\newcommand{\A}{\mathbb{A}}
\newcommand{\PP}{\mathbb{P}}
\newcommand{\E}{\mathbb{E}}
\newcommand{\eps}{\varepsilon}
\newcommand{\dd}{\,\mathrm{d}}
\begin{document}
\section*{UnsolvedMath ID: 30001985}
\textbf{Source problem code:} catalogue entry
\subsection{Definitions and mathematical statement}
Use the standard convention that a strongly minimal theory is a complete first-order theory \(T\) with infinite models such that, in every model \(M\models T\), every subset of its universe definable by a formula in one free variable, with parameters from \(M\), is finite or cofinite. Let \(L\) be the finite language of \(T\), with an effective coding of its symbols and formulas.

A countable model \(M\models T\) is recursively presentable if it is isomorphic to an \(L\)-structure with domain \(\N\) whose atomic diagram is decidable: there is an algorithm deciding every atomic or negated atomic sentence with named elements of the domain. Define the collections of isomorphism types
\[
\operatorname{Mod}_{\aleph_0}(T)=\{[M]_{\cong}:M\models T,\ |M|=\aleph_0\},
\qquad
\operatorname{Rec}(T)=\{[M]_{\cong}\in\operatorname{Mod}_{\aleph_0}(T):M\text{ is recursively presentable}\}.
\]
Let \(P_T\) denote the prime model of \(T\), meaning a model admitting an elementary embedding into every model of \(T\).

The source asks whether, for every strongly minimal theory \(T\) in a finite language satisfying the source's qualification ``Zilber's conjecture,'' one has
\[
\operatorname{Rec}(T)\in
\left\{\varnothing,\ \operatorname{Mod}_{\aleph_0}(T),\ \{[P_T]_{\cong}\}\right\}.
\]
Thus the third alternative means that exactly the isomorphism type of the prime model is recursively presentable. The supplied record does not state which formulation of Zilber's conjecture is intended, or precisely what it means for \(T\) to satisfy it. Accordingly, this is a conditional statement template retaining an unresolved source hypothesis, rather than a fully specified assertion about a determined class of theories.
\subsection{Short English statement}
For a strongly minimal theory in a finite language satisfying the unspecified Zilber hypothesis, must the recursively presentable countable models comprise no isomorphism types, every countable isomorphism type, or exactly the prime model's isomorphism type? The intended Zilber hypothesis requires clarification.
\subsection{Sources}
\begin{itemize}
\item Supplied local UnsolvedMath registry record: permanent ID \(30001985\), title \emph{Recursive Spectra of Strongly Minimal Theories}, source problem code \texttt{catalogue entry}, in the Oberwolfach Reports Open Problems collection. The record gives the proposed year as 2012 and names no proposer.
\item Supplied report citation: \emph{Computability Theory} (2012), Oberwolfach Reports, pp. 397--455, DOI \texttt{10.4171/owr/2012/08}. Source URL and supplied literature-source link, labeled ``original Oberwolfach report'': \url{https://doi.org/10.4171/owr/2012/08}. The background describes these reports as documenting mathematical workshops and their open-problem discussions.
\item The record labels the statement \texttt{corrected\_verified} and reports that its mathematical content was checked against the official report, with notation, context, and question form normalized without changing the source-level claim. These are supplied verification claims, not an independent verification in this draft.
\item Literature assessment dated exactly \texttt{2026-08-21}: the supplied curation labels the problem open and says that no primary source located through its search resolves the general conjecture. The only listed checked source is the original Oberwolfach report linked above. This is supplied dated literature triage, not independently verified; its time-sensitive status and cited claims require independent checking.
\end{itemize}
\end{document}
